
\section{\textbf{\texttt{grouped()}} ORM Method}

\begin{frame}[fragile]{\textbf{\texttt{grouped()}} Method}
	\begin{itemize}
		\item The \textbf{\texttt{grouped()}} method groups an existing recordset by a field or key.
		\item It is available on recordsets in modern Odoo versions (including Odoo 17).
		\item It returns a dictionary mapping each group key to a recordset.
\begin{block}{Method Syntax}
\begin{lstlisting}
def grouped(self, key):
	return super().grouped(key)
\end{lstlisting}
	\begin{itemize}
		\item Important parts: \textbf{\texttt{self}}, \textbf{\texttt{key}}.
	\end{itemize}
\end{block}
	\end{itemize}
\end{frame}

\begin{frame}[fragile]{Breaking Down the Method Signature}
	\begin{block}{}
		\begin{itemize}
			\item \textbf{\texttt{self}}: The recordset to group.
			\item \textbf{\texttt{key}}: A field name string, or a function that returns the group key.
\begin{lstlisting}
students.grouped('gender')
students.grouped('classroom_id')
students.grouped(lambda s: s.age // 10)  # age decade groups
\end{lstlisting}
		\end{itemize}
	\end{block}
\end{frame}

\begin{frame}[fragile]{What Does \texttt{grouped()} Return?}
	\begin{block}{}
		\begin{itemize}
			\item It returns a dictionary.
\begin{lstlisting}
students = self.env['student.student'].search([])
groups = students.grouped('gender')
print(groups)
\end{lstlisting}
			Output is (example):
\begin{lstlisting}
{
	'male': student.student(15, 17),
	'female': student.student(16, 22),
}
\end{lstlisting}
		\end{itemize}
	\end{block}
\end{frame}

\begin{frame}[fragile]{Grouping by a Many2one Field}
\begin{lstlisting}
grouped_by_class = students.grouped('classroom_id')

for classroom, class_students in grouped_by_class.items():
	print(classroom.name, len(class_students))
\end{lstlisting}
	\begin{itemize}
		\item Keys are relational records (or empty record / falsy key when empty).
		\item Values are recordsets of students in that group.
	\end{itemize}
\end{frame}

\begin{frame}[fragile]{Practical Use Cases}
	\begin{itemize}
		\item Build totals per group without writing SQL \texttt{GROUP BY}.
		\item Split records before batch processing.
		\item Prepare data for reports and dashboards.
\begin{lstlisting}
groups = students.grouped('gender')
male_count = len(groups.get('male', self.env['student.student']))
female_count = len(groups.get('female', self.env['student.student']))
\end{lstlisting}
	\end{itemize}
\end{frame}

\begin{frame}[fragile]{Method Call}
\begin{block}{Method Call}
\begin{lstlisting}
groups = students.grouped('classroom_id')
\end{lstlisting}
	\begin{itemize}
		\item Here, \textbf{\texttt{.grouped(...)}} is the \textbf{method call}.
		\item Return type reminder: \textbf{dict[key $\rightarrow$ recordset]}.
	\end{itemize}
\end{block}
\end{frame}
